\subsection{Shapley-Shubik Power Index}

\content{
        \begin{definition} A sequential coalition is an ordered coalition, that
        is, it keeps track of the order in which each player joins the
        coalition. 
        \end{definition}
}{% presentation
}{% nopresentation
}{% comments
We will use the notation 
$$ \la P_1, P_2, P_3\ra$$
to denote sequential coalitions. 
}

\content{
        \begin{definition} A \textbf{pivotal player} is the player in the
        coalition which changes the coalition from a losing coalition to a
        winning one.
        \end{definition}
}{% presentation
}{% nopresentation
}{% comments
}

\content{
        \begin{definition} (Shapley-Shubik power index) To compute the
        shapley-shubik power index: 
        \begin{itemize}
        \item List all sequential coalitions
        \item Determine all pivotal players
        \item Count up how many times each player was critical
        \item convert these into fractions by deviding each count by the total
        number of sequential coalitions
        \end{itemize}
        \end{definition}
}{% presentation
}{% nopresentation
}{% comments
}

\content{
        \begin{example} Calculate the Shapley-Shubik Power Index for the following
        weighted voting system
        $$ [ 8: 6,3,2 ]. $$
        \end{example}
}{% presentation
\vspace{3in}
}{% nopresentation
\vspace{3in}
}{% comments
}


\content{
        \begin{example} Compute the Shapley-Shubik Power Index for the weighted
        voting system: 
        $$ [ 36 : 20, 17, 15 ] $$
        \end{example}
}{% presentation
\vspace{3in}
}{% nopresentation
\vspace{3in}
}{% comments
$P_1 = 1/2$, $P_2 = 1/2$, $P_3 = 0$.
}

\content{
        \begin{example} Compute the Shapley-Shubik Power Index for the weighted
        voting system: 
        $$ [ 3: 2, 0.5, 0.5, 0.5 ] $$
        \end{example}
}{% presentation
\vspace{3in}
}{% nopresentation
\vspace{3in}
}{% comments
}

\content{
        \begin{example} Consider the weighted voting system given by 
        $$ [ q: 7, 3, 1 ]. $$
        Which values of $q$ result in a dictator? (list all values)
        \end{example}
}{% presentation
\vspace{3in}
}{% nopresentation
\vspace{3in}
}{% comments
}

\content{
        \begin{example} In the Shapley-Shubik method, is it possible for a dummy
        to be a pivital player? Why or why not? 
        \end{example}
}{% presentation
\vspace{2in}
}{% nopresentation
\vspace{2in}
}{% comments
No, if they were, they would be a critical player in a winning coalition. 
}
